\textbf{Phần 1. Trắc nghiệm lựa chọn}
\setcounterpageref{ex}{0}

\begin{ex} %[4D3N13]
Khi làm tròn số đến hàng trăm nghìn, ta so sánh chữ số ở hàng nào với 5?
\choice
{Hàng trăm nghìn}
{\True Hàng chục nghìn}
{Hàng nghìn}
{Hàng trăm}
\loigiai{
Ta so sánh chữ số ở hàng chục nghìn với 5.\\
\textbf{Đáp án của câu hỏi là B}
}
\end{ex}

\begin{ex} %[4D3N13]
Làm tròn số $430\,000$ đến hàng trăm nghìn ta được
\choice
{\True $400\,000$}
{$500\,000$}
{$450\,000$}
{$430\,000$}
\loigiai{
Chữ số hàng chục nghìn là $3 < 5$ nên làm tròn thành $400\,000$.\\
\textbf{Đáp án của câu hỏi là A}
}
\end{ex}

\begin{ex} %[4D3H13]
Làm tròn số $670\,000$ đến hàng trăm nghìn ta được
\choice
{$600\,000$}
{\True $700\,000$}
{$650\,000$}
{$800\,000$}
\loigiai{
Chữ số hàng chục nghìn là $7 \ge 5$ nên làm tròn thành $700\,000$.\\
\textbf{Đáp án của câu hỏi là B}
}
\end{ex}

\begin{ex} %[4D3H13]
Làm tròn số $1\,250\,000$ đến hàng trăm nghìn ta được
\choice
{$1\,200\,000$}
{$1\,000\,000$}
{\True $1\,300\,000$}
{$2\,000\,000$}
\loigiai{
Chữ số hàng chục nghìn là 5 nên làm tròn lên $1\,300\,000$.\\
\textbf{Đáp án của câu hỏi là C}
}
\end{ex}

\begin{ex} %[4D3H13]
Số nào sau đây khi làm tròn đến hàng trăm nghìn thì được $500\,000$?
\choice
{$435\,000$}
{$560\,000$}
{$390\,000$}
{\True $482\,000$}
\loigiai{
$482\,000$ có chữ số hàng chục nghìn là $8 \ge 5$ nên làm tròn thành $500\,000$.\\
\textbf{Đáp án của câu hỏi là D}
}
\end{ex}

\begin{ex} %[4D3H13]
Làm tròn số $815\,420$ đến hàng trăm nghìn ta được
\choice
{$900\,000$}
{$820\,000$}
{\True $800\,000$}
{$850\,000$}
\loigiai{
Chữ số hàng chục nghìn là $1 < 5$ nên làm tròn thành $800\,000$.\\
\textbf{Đáp án của câu hỏi là C}
}
\end{ex}

\begin{ex} %[4D3V13]
Dân số một thành phố là $2\,845\,600$ người. Làm tròn dân số đến hàng trăm nghìn ta được
\choice
{\True $2\,800\,000$ người}
{$2\,900\,000$ người}
{$3\,000\,000$ người}
{$2\,850\,000$ người}
\loigiai{
Hàng chục nghìn là $4 < 5$ nên làm tròn thành $2\,800\,000$ người.\\
\textbf{Đáp án của câu hỏi là A}
}
\end{ex}

\begin{ex} %[4D3V13]
Một công ty xuất khẩu được $3\,652\,000$ tấn gạo. Làm tròn số tấn gạo đến hàng trăm nghìn ta được
\choice
{$3\,600\,000$ tấn}
{$4\,000\,000$ tấn}
{$3\,650\,000$ tấn}
{\True $3\,700\,000$ tấn}
\loigiai{
Hàng chục nghìn là $5 \ge 5$ nên làm tròn lên $3\,700\,000$ tấn.\\
\textbf{Đáp án của câu hỏi là D}
}
\end{ex}

\begin{ex} %[4D3V13]
Tìm số tự nhiên lớn nhất có 6 chữ số mà khi làm tròn đến hàng trăm nghìn ta được $600\,000$.
\choice
{$649\,000$}
{$650\,000$}
{\True $649\,999$}
{$599\,999$}
\loigiai{
Số lớn nhất là $649\,999$ (hàng chục nghìn là 4).\\
\textbf{Đáp án của câu hỏi là C}
}
\end{ex}

\begin{ex} %[4D3V13]
Tìm số tự nhiên bé nhất có 6 chữ số mà khi làm tròn đến hàng trăm nghìn ta được $600\,000$.
\choice
{$549\,999$}
{\True $550\,000$}
{$500\,000$}
{$555\,555$}
\loigiai{
Số bé nhất là $550\,000$ (hàng chục nghìn là 5).\\
\textbf{Đáp án của câu hỏi là B}
}
\end{ex}

\begin{ex} %[4D3V13]
Làm tròn số $950\,000$ đến hàng trăm nghìn ta được
\choice
{$900\,000$}
{$950\,000$}
{$1\,100\,000$}
{\True $1\,000\,000$}
\loigiai{
Chữ số hàng chục nghìn là 5 nên $950\,000$ làm tròn thành $1\,000\,000$.\\
\textbf{Đáp án của câu hỏi là D}
}
\end{ex}

\begin{ex} %[4D3C13]
Có bao nhiêu số tròn chục nghìn khi làm tròn đến hàng trăm nghìn thì được số $800\,000$?
\choice
{\True $10$ số}
{$9$ số}
{$11$ số}
{$5$ số}
\loigiai{
Các số đó là: $750\,000, 760\,000, \dots, 790\,000$ (5 số) và $800\,000, 810\,000, \dots, 840\,000$ (5 số). Tổng cộng có $5 + 5 = 10$ số.\\
\textbf{Đáp án của câu hỏi là A}
}
\end{ex}

\textbf{Phần 2. Trắc nghiệm đúng sai}
\setcounterpageref{ex}{0}

\begin{ex} %[4D3H13]
Xét tính đúng sai của việc làm tròn số đến hàng trăm nghìn:
\choiceTF
{\True $348\,200$ làm tròn đến hàng trăm nghìn được $300\,000$.}
{\True $762\,500$ làm tròn đến hàng trăm nghìn được $800\,000$.}
{$450\,000$ làm tròn đến hàng trăm nghìn được $400\,000$.}
{\True $450\,000$ làm tròn đến hàng trăm nghìn được $500\,000$.}
\loigiai{
\begin{itemize}
    \item $348\,200$: chục nghìn là 4 nên làm tròn $300\,000$. Đúng.
    \item $762\,500$: chục nghìn là 6 nên làm tròn $800\,000$. Đúng.
    \item $450\,000$: chục nghìn là 5 nên làm tròn lên $500\,000 \neq 400\,000$. Ý 3 Sai, Ý 4 Đúng.
\end{itemize}
\begin{itemize}
    \item a) \textbf{Đúng}.
    \item b) \textbf{Đúng}.
    \item c) \textbf{Sai}.
    \item d) \textbf{Đúng}.
\end{itemize}
}
\end{ex}

\begin{ex} %[4D3V13]
Cho các số: $512\,000; 549\,900; 580\,000; 635\,000$. Xét tính đúng sai:
\choiceTF
{\True Có 2 số khi làm tròn đến hàng trăm nghìn được $500\,000$ ($512\,000$ và $549\,900$).}
{\True Có 2 số khi làm tròn đến hàng trăm nghìn được $600\,000$ ($580\,000$ và $635\,000$).}
{Số $580\,000$ làm tròn đến hàng trăm nghìn được $500\,000$.}
{\True Làm tròn số $635\,000$ đến hàng trăm nghìn được $600\,000$.}
\loigiai{
\begin{itemize}
    \item $512\,000$ và $549\,900 \rightarrow 500\,000$. Đúng.
    \item $580\,000$ và $635\,000 \rightarrow 600\,000$. Đúng.
    \item $580\,000 \rightarrow 600\,000 \neq 500\,000$. Sai.
    \item $635\,000 \rightarrow 600\,000$. Đúng.
\end{itemize}
\begin{itemize}
    \item a) \textbf{Đúng}.
    \item b) \textbf{Đúng}.
    \item c) \textbf{Sai}.
    \item d) \textbf{Đúng}.
\end{itemize}
}
\end{ex}

\textbf{Phần 3. Trả lời ngắn}
\setcounterpageref{ex}{0}

\begin{ex} %[4D3H13]
Làm tròn số $724\,800$ đến hàng trăm nghìn.
\shortans{700000}
\loigiai{
\textbf{Đáp số:} $700000$\\
Hàng chục nghìn là 2 nên làm tròn thành $700\,000$.
}
\end{ex}

\begin{ex} %[4D3H13]
Làm tròn số $876\,500$ đến hàng trăm nghìn.
\shortans{900000}
\loigiai{
\textbf{Đáp số:} $900000$\\
Hàng chục nghìn là 7 nên làm tròn lên $900\,000$.
}
\end{ex}

\begin{ex} %[4D3V13]
Làm tròn số $2\,450\,000$ đến hàng trăm nghìn.
\shortans{2500000}
\loigiai{
\textbf{Đáp số:} $2500000$\\
Hàng chục nghìn là 5 nên làm tròn lên $2\,500\,000$.
}
\end{ex}

\begin{ex} %[4D3V13]
Tìm số tự nhiên lớn nhất có 6 chữ số khác nhau mà khi làm tròn đến hàng trăm nghìn được $400\,000$.
\shortans{449876}
\loigiai{
\textbf{Đáp số:} $449876$\\
$449\,876$ (chữ số hàng chục nghìn là 4).
}
\end{ex}

\begin{ex} %[4D3V13]
Tìm số tự nhiên bé nhất có 6 chữ số khác nhau mà khi làm tròn đến hàng trăm nghìn được $400\,000$.
\shortans{350124}
\loigiai{
\textbf{Đáp số:} $350124$\\
$350\,124$ (chữ số hàng trăm nghìn là 3, chục nghìn là 5).
}
\end{ex}

\begin{ex} %[4D3C13]
Một huyện có $184\,600$ người, sau một năm tăng thêm $18\,400$ người. Làm tròn tổng dân số sau một năm đến hàng trăm nghìn.
\shortans{200000}
\loigiai{
\textbf{Đáp số:} $200000$\\
Tổng dân số: $184\,600 + 18\,400 = 203\,000 \rightarrow 200\,000$ người.
}
\end{ex}
